\section{模型假设}

\textbf{假设一}：电解槽与合成氨装置作为整体同步运行，即制氢设备与合成氨装置同时开机或停机，不存在单独运行某一设备的情况。该假设基于题目明确的"全额开机和停机两种运行方式"描述，且制氢直接供合成氨的工艺流程要求两者必须同步。

\textbf{假设二}：设备功率随产能规模呈线性同步提升，扩容倍数$k = Q_{\text{NH3}} / 36$。该假设直接来源于题目条件"配套电氢氨装置的额定功率将随产能规模呈线性同步提升"，反映了工业园区的模块化扩建特征。

\textbf{假设三}：不计园区内部功率损耗，即忽略变压器损耗、线路损耗和电力电子变换损耗。该假设由题目问题一明确给出，在实际工程中园区内部损耗通常不超过2\%--3\%，对整体分析影响有限。

\textbf{假设四}：同一时段不能同时购电和售电，即$P_{\text{buy}}(t) \cdot P_{\text{sell}}(t) = 0$。该假设基于物理约束——园区与电网仅有一个并网接口，功率方向在任一时刻唯一确定。

\textbf{假设五}：风光发电的度电成本（风电0.15元/kWh、光伏0.12元/kWh）已包含设备投资折旧和运维费用，无需另行计算风光设备的资本折旧。合成氨装置投资成本按30年使用寿命计算日均折旧。该假设基于附件5给出的度电成本定义方式。
